The development of large language models (LLMs) has afforded tremendous advancements in natural language generation (NLG). 
Recently, LLMs have been widely applied across various natural language domains \citep{zhang2023sentiment,wang2023gpt,alves-etal-2023-steering,zhang-etal-2024-la-ucl}, even extending to tasks and domains traditionally dominated by other machine learning algorithms, such as graph data~\citep{fatemi2024talk}, tabular data~\citep{borisov2023language, hegselmann2023tabllm}, time series~\citep{gruver2023llmtime, rasul2024lagllama}, predictive chemistry~\citep{Jablonka2024LeveragingLL, shi2023relm}, computer vision~\citep{wang2023visionllm}, amongst others.
As LLMs continue to demonstrate outstanding performance, their reliability is increasingly scrutinized. 
Uncertainty quantification, an area of research that can provide some guidance on reliability, has recently gained much attention. 
Despite its long history in other machine learning tasks \citep{gawlikowski2023survey}, our understanding of uncertainty quantification in NLG remains relatively limited. %, largely due to several challenges unique to NLG.

During pre-training, (autoregressive) LLMs are optimized to predict high logits for the target token to minimize (variants of) the negative log likelihood. 
Consequently, one of the most natural and widely used confidence scores in selective NLG, conformal NLG, or uncertainty quantification is the (sometimes normalized) likelihood of the sequence, equivalent to the sum/mean of token logits. 
At first glance, sequence likelihood appears to be the most faithful reflection of a model's confidence, as it represents the (log of) the predicted probability of the output sequence, or $\log{\hat{p}(\predSeq|x)}$. 
However, this measure often lacks proper contextualization and disregards the specific nature of the task at hand. 
$\hat{p}(\predSeq|x)$ conflates syntactic and semantic likelihoods, though we typically prioritize the semantic aspect, to a varying degree depending on the task. 
For example, depending on whether the question is ``Which country won the World Cup in 2022'' or ``When did Messi win the World Cup,'' the answer ``Messi emerged victorious in the 2022 World Cup'' could be considered correct or incorrect. 
Further, the somewhat unusual expression here could adversely impact $\hat{p}(\predSeq|x)$.

Despite its limitations, relatively limited attention has been paid to improving the vanilla sequence likelihood. 
A recent study by ~\citet{duan2023shifting} proposed assessing token relevance using an external natural language inference (NLI) model. 
The relevance score of a token is negatively proportional to the similarity between the original sequence and the sequence with this token removed, which is then applied to weight the token logits.
%This token relevance score is calculated by removing each token from the original generated sequence to create a new sequence, measuring the similarity between the new sequence and the original, and normalizing these similarity weights so that they sum to one. 
%The normalized weight is then applied to the token logits.
However, since modern LMs rely on sub-word tokenizers, removing one sub-word token at a time results in non-words, is computationally expensive for longer texts, and remains context-unaware.


%while the sequence logits reflect the likelihood of $\predSeq$ given prompt $x$, it ignores the nature of the prompt

In this paper, we investigate the potential of utilizing the LLM's own attention mechanism to develop a weighted sequence likelihood as an enhanced confidence score. 
Specifically, the LLM is prompted to concentrate on the relevant tokens in its own generation. 
The most appropriate attention heads are then employed to re-weight the original token logits. The main contributions of this paper are summarized as thus:
\begin{itemize}

\item We introduce a straightforward yet effective method to reweight token logits, resulting in \underline{C}ontextualized \underline{S}equence \underline{L}ikelihood (\methodname), a more reliable confidence measure.
        
\item We improve current automatic evaluation methods for confidence measures on Question-Answering (QA) datasets and manually verify their effectiveness.

\item In popular free-form QA datasets, and on a variety of LLMs, we verify that \methodname significantly outperforms baselines.
Case studies suggest the attention weights are meaningful.
\end{itemize}


